\PassOptionsToPackage{dvipsnames}{xcolor}
\documentclass[portrait,a1paper,fontscale=0.45,]{baposter}

\usepackage{capt-of}
\usepackage{graphicx}
\usepackage{multicol}
\usepackage{tikz}

\graphicspath{{../../../logos}, {figures/article2},
  {../../../articles/extended-abstract-article2/figures},{figures/rules}}

\selectcolormodel{HTML}
\definecolor{pastelpink}{RGB}{255,192,255}
\definecolor{warmorange}{RGB}{255,128,0}
\definecolor{myorange}{RGB}{255,150,0}
\definecolor{myorange2}{RGB}{255,192,0}
\definecolor{mygreen}{RGB}{0,192,0}
\definecolor{mygreen2}{RGB}{0,208,0}


\begin{document}
\begin{poster}%
  {%
    grid=false,%
    eyecatcher=true,%
    columns=6,%
    background=none,%
    % headerborder=none,%
    headerborder=none,%
    headerfont=\LARGE\bfseries\sffamily,%
    headerFontColor=white,%
    headershade=plain,%
    textborder=none,%
    % borderColor=red,%
    boxshade=none}%
  {%
    \begin{minipage}[c]{9em}
      \center
    
\includegraphics[width=.6\textwidth]{branding_eguk.gif.png}\\
    \vspace{1em}
    
\includegraphics[width=\textwidth]{cgvc-logo}
    \vspace{1em}
    2023
    \end{minipage}
  }%
  {Model reevaluation based on graph transformation rules
  \vspace{.2em}}%
{M. Gaide, D. Marcheix, A. Arnould, X. Skapin, H. Belhaouari \& S. Jean\\
  maxime.gaide@ensma.fr
  }%
  {%
    \begin{minipage}[c]{14em}
      
\includegraphics[width=6em]{Logo_ISAE-ENSMA}
      \hfill
      \includegraphics[width=6em]{Logo_up}

      \vspace{1em}

      
\includegraphics[width=6em]{lias-logo}
      \hfill
      
\includegraphics[width=6em]{xlim-logo}
    \end{minipage}
  }

  \begin{posterbox}[name=abstract,row=0, span=3]{Abstract}
    % \begin{minipage}[c]{0.89\linewidth}
    We extend the scope of naming problem studies to encompass rule-based graph
    transformation modeling systems. We propose a new persistent naming method
    that capitalizes on the formalized operations of generalized maps and graph
    transformation rules. It enables a unique and homogeneous characterisation
    of entities across all dimensions.
    % \end{minipage}
    % \hfill
    % \begin{minipage}[c]{0.08\linewidth}
    %  \center
     % 
\includegraphics[width=\textwidth]{reevaluation-demo.png}
     % Demo
    % \end{minipage}
  \end{posterbox}

  \begin{posterbox}[name=context, column=3, span=3]{Context}
    \underline{\bfseries Evaluation}

      \includegraphics[width=.19\textwidth]{fil-rouge-short-create-face-magnified-eval}
      \includegraphics[width=.19\textwidth]{fil-rouge-short-extrude-face-magnified}
      \includegraphics[width=.19\textwidth]{fil-rouge-short-triangulate-face-magnified}
      \includegraphics[width=.195\textwidth]{fil-rouge-short-color-face-magnified}

      1-create(pos); 2-extrude($PN_1$, vec); 3-triangulate($PN_2$); 4-color($PN_3$)

    \vspace{1em}

    \underline{\bfseries Reevaluation}

    \includegraphics[width=.19\textwidth]{fil-rouge-short-create-face-magnified-reval}
    \includegraphics[width=.19\textwidth]{fil-rouge-short-reeval-insert-vertex-magnified}
    \includegraphics[width=.19\textwidth]{fil-rouge-short-reeval-extrude-face-magnified}
    \includegraphics[width=.19\textwidth]{fil-rouge-short-reeval-triangulate-face-magnified}
    \includegraphics[width=.19\textwidth]{fil-rouge-short-reeval-color-face-magnified}

    1-create(pos); ADD-insert(3); 2-extrude($PN_1$, vec); 3-triangulate($PN_2$);\\ 4-color($PN_3$)
  \end{posterbox}

  \begin{posterbox}[name=rules, below=abstract, column=0, span=3]{Rules}
    \begin{minipage}[c]{\linewidth}
      \center
      \includegraphics[width=.6\textwidth]{rule-create-square-face-poster-version}\\
      Square creation rule

      \vspace{1em}

      \includegraphics[height=40pt]{rule-face-extrusion-poster-version}\\
      Face extrusion rule

      \vspace{1em}

      \includegraphics[height=35pt]{rule-face-triangulation-3D-poster-version}\\
      Face triangulation rule
    \end{minipage}

  \end{posterbox}


  \begin{posterbox}[name=persistent name, below=context, column=3, span=3]{Persistent naming scheme}
    Taking advantage of Jerboa graph transformation rules on generalized map's formalism:
    \begin{itemize}
      \item Using darts to represent orbits (cells)
      \item Persistent name: history of a darts' creation
    \end{itemize}
    Persistent name
    $PN_{3}=[1\fcolorbox{black}{magenta}{$n_{3}$};2\fcolorbox{black}{warmorange}{$n_{4}$};3\fcolorbox{black}{mygreen}{$n_{0}$}]$
    \begin{itemize}
      \item Rule \texttt{1-create}'s node \fcolorbox{black}{magenta}{$n_{3}$} creates dart \fcolorbox{red}{pastelpink}{3}
      \item Rule \texttt{2-extrude} matches dart \fcolorbox{red}{pastelpink}{3} and
            node \fcolorbox{black}{warmorange}{$n_{4}$} creates its copy: dart
            \fcolorbox{red}{myorange}{35}
      \item Rule \texttt{3-triangulate}'s node \fcolorbox{black}{mygreen}{$n_{0}$} matches dart \fcolorbox{red}{mygreen2}{35}
    \end{itemize}
  \end{posterbox}

  % \begin{posterbox}[name=rules two, headershade=shadelr,headerColorOne=white,headerColorTwo=white,below=persistent name, span=2]{}
  %   \begin{minipage}[c]{.5\linewidth}
  %     \center
  %     \includegraphics[height=40pt]{rule-face-extrusion-poster-version}\\
  %     Face extrusion rule
  %   \end{minipage}
  %   \begin{minipage}[c]{.5\linewidth}
  %     \center
  %     \includegraphics[height=35pt]{rule-face-triangulation-3D-poster-version}\\
  %     Face triangulation rule
  %   \end{minipage}
  % \end{posterbox}

  \begin{posterbox}[name=evaluation, below=persistent name, column=0, span=3]{Evaluation
      DAG of PN$_{3}$}
    \begin{itemize}
      \item Represents a topological parameter's history
      \item Built from a persistent name
    \end{itemize}
      \includegraphics[height=25em]{pn7-short}
  \end{posterbox}

  \begin{posterbox}[name=reevaluation, below=persistent name, column=3,
    span=3]{Reevaluation DAG of PN$_{3}$}
    \begin{itemize}
      \item Represents the matched topological parameter's histories
      \item Built from an evaluation DAG
    \end{itemize}
    \includegraphics[height=25em]{pn7-reval-short}
  \end{posterbox}

  \begin{posterbox}[name=conclusion, below=evaluation, column=0, span=5]{Conclusion}
    \setlength{\columnseprule}{1pt}
    \def\columnseprulecolor{\color{gray}}

    \begin{multicols}{2}
      \begin{minipage}[c]{\linewidth}
        \begin{itemize}
          \item Rule-based persistent naming scheme
          \begin{itemize}
            \item Persistent naming scheme based on rules
            \item Unique identification of topological parameters through their histories
          \end{itemize}
          \item Add, Delete, Move operations to edit a model.
          \end{itemize}
        \end{minipage}

      \columnbreak

      \begin{minipage}[c]{\linewidth}
        \begin{itemize}
          \item Method:
          \begin{itemize}
            \item Record darts histories
            \item Rebuild orbit histories (evaluation DAG)
            \item Build new orbit histories (reevaluation DAG) using strategies
          \end{itemize}
        \end{itemize}
      \end{minipage}
    \end{multicols}
  \end{posterbox}

  \begin{posterbox}[name=Demo, below=evaluation, column=5]{Demo}
    \center
    \vspace{1em}
    
\includegraphics[scale=0.08]{reevaluation-demo.png}\\
  \end{posterbox}

  % \begin{posterbox}[name=link, below=evaluation, column=0]{Demo}
  %   An overview of this work's implementation is available through this link
  %   \includegraphics[width=.15\textwidth]{qr-code-reeval}\
  % \end{posterbox}

  % \begin{posterbox}[name=participants, below=link]{Participants}
  %   \begin{tabular}{l l}
  %     \textbf{Maxime Gaide} & \textbf{Xavier Skapin}\\
  %     \small{maxime.gaide@ensma.fr} & \small{xavier.skapin@univ-poitiers.fr}\\
  %     \textbf{David Marcheix} & \textbf{Hakim Belhaouari}\\
  %     \small{david.marcheix@ensma.fr} & \small{hakim.ferrier-belhaouari@univ-poitiers.fr}\\
  %     \textbf{Agnès Arnould} & \textbf{Stéphane Jean}\\
  %     \small{agnes.arnould@univ-poitiers.fr} & {stephane.jean@ensma.fr}
  %   \end{tabular}
  % \end{posterbox}

  \begin{posterbox}[name=CGVC, above=bottom, span=6, textborder=none,
    headerborder=none, textfont=\Large\bf\sffamily]{\centerline{Computer Graphics
      and Visual Computing 2023}}
  \end{posterbox}

\end{poster}
\end{document}
